\section{RELATED WORK AND BACKGROUND}

\subsection{Edge-Server Split Inference}

Edge-server collaborative inference distributes computation across edge and server
resources to avoid executing full MLLM pipelines on constrained devices. Prior work has
demonstrated, however, that na\"{i}ve layer-wise splitting increases end-to-end latency
by $2.35\times$ compared to server-only execution due to the per-token communication
overhead inherent in autoregressive generation~\cite{park2025specedge}. Approaches such
as EdgeShard~\cite{zhang2025edgeshard}, Splitwise~\cite{younesi2025splitwise}, and
adaptive layer splitting~\cite{chen2025adaptivesplit} address this through collaborative
layer partitioning and dynamic cut-point selection, reducing communication overhead while
balancing computation across the edge-server boundary. However, these systems are designed
exclusively for single-modality text pipelines and do not consider the additional challenge
of encoding and aligning heterogeneous sensor streams before transmission.
{\em EMMA addresses this gap by performing lightweight multimodal encoding and cross-modal
alignment at the edge and transmitting a single fixed-size latent vector, avoiding
per-token cross-boundary communication regardless of the number of input modalities.}

\subsection{On-Device Optimization}

A parallel line of work reduces MLLM resource requirements to enable on-device deployment.
Dedicated accelerators such as NPUs, FPGAs, and NPU-PIM hybrids~\cite{p3llm,motoyama2024iedm}
target the memory-bandwidth bottleneck of transformer inference through near-data compute,
while activation-aware quantization methods such as AWQ~\cite{lin2024awq} enable 4-bit
mixed-precision execution on commodity edge GPUs with negligible accuracy loss. Knowledge
distillation approaches~\cite{yao2024gkt,yao2025efficient} offer a complementary route,
producing compact edge-deployable MLLMs by transferring capability from large server-side
models. Within single-device MLLM inference, visual token reduction methods further reduce
computation: FastV~\cite{chen2024fastv} prunes visual tokens in deep transformer layers
based on attention patterns, LLaVA-PruMerge~\cite{shang2024llavaprumerge} performs adaptive
token selection at the feature projection stage, and VisionZip~\cite{yang2024visionzip}
selects compact visual representations by identifying redundant tokens through encoder
attention analysis. Despite these advances, the fundamental capacity gap between modern
MLLMs and fully constrained edge platforms remains: processing a single image-text query
with LLaVA-1.5 requires approximately 18.2 TFLOPs and 41.6 GB of memory, far exceeding
the 1.3 TFLOPs and 8 GB available on platforms such as the NVIDIA Jetson TX2. Furthermore,
as noted in~\cite{yuan2025task}, visual token reduction methods assume inference is executed
on a single device and do not consider transmission constraints in edge-server pipelines.
{\em On-device optimization alone cannot support full MLLM inference on the most constrained
edge devices, motivating EMMA's approach of offloading reasoning to the server while keeping
only lightweight encoding at the edge.}

\subsection{Feature Compression and Offloading}

A growing body of work applies task-oriented communication principles to reduce transmission
overhead in edge inference pipelines, compressing features to retain only task-relevant
information rather than optimizing for pixel-accurate reconstruction.

He et al.~\cite{he2022elic} propose ELIC, an efficient learned image codec using
space-channel contextual adaptive entropy coding. ELIC optimizes for pixel-accurate
reconstruction and serves as a general-purpose image transmission baseline, but has no
awareness of downstream task semantics — compression quality is measured in PSNR/MS-SSIM,
not task accuracy.

Yuan et al.~\cite{yuan2025task} propose a task-oriented feature compression (TOFC) method
for large multimodal model inference in a device-edge co-inference framework. TOFC employs
density-peaks clustering to merge visual features, reducing the number of transmitted visual
tokens to only 1.1\%--4.4\% of the original count, and applies a learnable entropy model
with hyperprior to further compress the merged features. Crucially, TOFC treats modalities
independently: the textual instruction is transmitted raw to the server, and cross-modal
alignment occurs entirely on the server side, meaning image compression at the edge is not
informed by the user query. Furthermore, TOFC trains the compressor end-to-end with the
downstream task model, requiring task labels and a fixed server model during compression
training. MoA-Off~\cite{yang2025moaoff} introduces a lightweight modality-aware module that
estimates the complexity of heterogeneous inputs and dynamically schedules which modalities
are processed locally versus offloaded to the cloud, achieving over 30\% latency reduction
and 30\%--65\% decrease in resource overhead. Semantic communication
frameworks~\cite{ref19,ref20} transmit task-relevant representations rather than raw data,
providing theoretical grounding for learned compact encodings at the edge, but are
predominantly designed for single-modality streams and require end-to-end joint optimization
of encoder and decoder.

To compress the latent itself, we adapt the Relational Knowledge Distillation (RKD)
objective of Park et al.~\cite{park2019relational}, which preserves pairwise structural
relationships between samples rather than matching individual outputs. Rather than distilling
from teacher to student model, EMMA distills geometric structure from a high-dimensional
CLIP embedding into a bandwidth-constrained latent vector — the DistilAE loss — without
requiring task labels.

{\em EMMA differs from these works in three key respects. First, EMMA performs cross-modal
alignment at the edge prior to compression, transmitting a single unified $k$-dimensional
latent regardless of input complexity. Second, EMMA's compressor is trained without task
labels — supervision comes entirely from the pre-trained CLIP embedding geometry via the
DistilAE relational distillation loss — enabling the edge compressor to be deployed
independently of any specific server task. Third, no decoder runs at inference: the
compressed latent is mapped directly to soft prompt tokens via a lightweight projection MLP
($\sim$100K parameters), eliminating server-side reconstruction overhead and allowing the
frozen server LLM to be upgraded or replaced without retraining the edge components.}
